\documentclass[a4paper,10pt]{report}\input{../0.Préambule/Préambule_cours_prof}\begin{document}

\pagestyle{empty}
\newpage
\noindent NOM, Prénom et classe : \dotfill

\section*{Triangles et quadrilatères}
\addcontentsline{toc}{section}{Triangles et quadrilatères}

\noindent \textit{L'usage de la calculatrice est \textbf{interdite}.} \hfill \textit{Durée : 55 minutes}


\medskip
\paragraph{Exercice 1}: \hfill \textbf{3 points}

\noindent \begin{minipage}{5.5cm}
\begin{center}
\begin{tikzpicture}[line cap=round,line join=round,>=triangle 45,x=1.0cm,y=1.0cm,scale=0.3]
\clip(0.5,2.) rectangle (9.5,11.);
\draw [line width=1.2pt,domain=0.5:9.5] plot(\x,{(--87.0016-1.7*\x)/8.12});
\draw [line width=1.2pt,domain=0.5:9.5] plot(\x,{(--33.0644-6.22*\x)/-2.78});
\draw [line width=1.2pt,domain=1:7] plot(\x,{(-13.6896--0.32*\x)/-4.54});
\draw [line width=1.2pt,domain=0.5:9.5] plot(\x,{(-16.896--7.6*\x)/-0.8});
\begin{scriptsize}
\draw [color=blue] (1.12,10.48)-- ++(-3.5pt,0 pt) -- ++(7.0pt,0 pt) ++(-3.5pt,-3.5pt) -- ++(0 pt,7.0pt);
\draw[color=blue] (1.84,9.7) node {\Large{$A$}};
\draw [color=blue] (9.24,8.78)-- ++(-3.5pt,0 pt) -- ++(7.0pt,0 pt) ++(-3.5pt,-3.5pt) -- ++(0 pt,7.0pt);
\draw[color=blue] (8.7,9.65) node {\Large{$B$}};
\draw [color=blue] (6.46,2.56)-- ++(-3.5pt,0 pt) -- ++(7.0pt,0 pt) ++(-3.5pt,-3.5pt) -- ++(0 pt,7.0pt);
\draw[color=blue] (6.16,3.27) node {\Large{$C$}};
\draw [color=blue] (1.92,2.88)-- ++(-3.5pt,0 pt) -- ++(7.0pt,0 pt) ++(-3.5pt,-3.5pt) -- ++(0 pt,7.0pt);
\draw[color=blue] (2.36,3.63) node {\Large{$D$}};
\end{scriptsize}
\end{tikzpicture}
\end{center}
\end{minipage}
\begin{minipage}{11.5cm}
\begin{enumerate}
\item Quel est le côté opposé au segment $[BC]$ dans le quadrilatère $ABCD$ ?

\medskip\noindent\grandscarreaux{\linewidth}{1.6cm}
\end{enumerate}
\end{minipage}

\begin{enumerate}
\setcounter{enumi}{1}
\item Quels sont les côtés consécutifs au segment $[AD]$ dans le quadrilatère $ABCD$ ?

\medskip\noindent\grandscarreaux{\linewidth}{1.6cm}
\item Quelles sont les diagonales du quadrilatère $ABCD$ ?

\medskip\noindent\grandscarreaux{\linewidth}{1.6cm}
\end{enumerate}

\medskip
\paragraph{Exercice 2}: \hfill \textbf{5 points}

\begin{center}
\definecolor{qqwuqq}{rgb}{0.,0.39215686274509803,0.}
\begin{tikzpicture}[line cap=round,line join=round,>=triangle 45,x=1.0cm,y=1.0cm,scale=0.8]
\clip(-3.,1.3) rectangle (10.,5.3);
\draw[line width=1.2pt,color=qqwuqq,fill=qqwuqq,fill opacity=0.1] (7.689339398899502,4.217700039357383) -- (8.031639359542119,3.9670394382568848) -- (8.282299960642616,4.3093393988995015) -- (7.94,4.56) -- cycle; 
\draw [line width=1.2pt] (3.,5.)-- (2.,2.);
\draw [line width=1.2pt] (2.6201665510863985,3.544271887242357) -- (2.4304298914762956,3.607517440445725);
\draw [line width=1.2pt] (2.5948683298050517,3.468377223398316) -- (2.405131670194949,3.531622776601684);
\draw [line width=1.2pt] (2.5695701085237044,3.392482559554275) -- (2.3798334489136015,3.4557281127576434);
\draw [line width=1.2pt] (5.,3.)-- (2.,2.);
\draw [line width=1.2pt] (3.6075174404457244,2.4304298914762956) -- (3.5442718872423575,2.6201665510863985);
\draw [line width=1.2pt] (3.5316227766016834,2.405131670194949) -- (3.4683772233983166,2.5948683298050517);
\draw [line width=1.2pt] (3.4557281127576425,2.3798334489136015) -- (3.3924825595542756,2.5695701085237044);
\draw [line width=1.2pt] (3.,5.)-- (5.,3.);
\draw [line width=1.2pt] (-2.6,2.44)-- (0.18,1.62);
\draw [line width=1.2pt] (-1.2200744774407992,2.137231081801383) -- (-1.276657158957616,1.9454019908053442);
\draw [line width=1.2pt] (-1.143342841042384,2.1145980091946557) -- (-1.199925522559201,1.9227689181986172);
\draw [line width=1.2pt] (0.18,1.62)-- (-0.49985916889675985,4.437550622520741);
\draw [line width=1.2pt] (-0.24775718664326518,2.966435016966094) -- (-0.0533369798974104,3.0133475228563054);
\draw [line width=1.2pt] (-0.26652218899934954,3.0442030996644363) -- (-0.07210198225349472,3.0911156055546476);
\draw [line width=1.2pt] (-0.49985916889675985,4.437550622520741)-- (-2.6,2.44);
\draw [line width=1.2pt] (-1.4520275048126956,3.3938845237532647) -- (-1.5898650300417336,3.538801108859092);
\draw [line width=1.2pt] (-1.5099941388550264,3.3387495136616496) -- (-1.6478316640840642,3.4836660987674772);
\draw [line width=1.2pt] (6.08,2.02)-- (7.94,4.56);
\draw [line width=1.2pt] (6.08,2.02)-- (9.51142979696883,3.409267943951958);
\draw [line width=1.2pt] (7.94,4.56)-- (9.51142979696883,3.409267943951958);
\begin{scriptsize}
\draw [color=blue] (3.,5.)-- ++(-3.5pt,0 pt) -- ++(7.0pt,0 pt) ++(-3.5pt,-3.5pt) -- ++(0 pt,7.0pt);
\draw[color=blue] (3.58,5.13) node {\large{$D$}};
\draw [color=blue] (2.,2.)-- ++(-3.5pt,0 pt) -- ++(7.0pt,0 pt) ++(-3.5pt,-3.5pt) -- ++(0 pt,7.0pt);
\draw[color=blue] (1.8,2.61) node {\large{$E$}};
\draw [color=blue] (5.,3.)-- ++(-3.5pt,0 pt) -- ++(7.0pt,0 pt) ++(-3.5pt,-3.5pt) -- ++(0 pt,7.0pt);
\draw[color=blue] (4.72,2.69) node {\large{$F$}};
\draw [color=blue] (-2.6,2.44)-- ++(-3.5pt,0 pt) -- ++(7.0pt,0 pt) ++(-3.5pt,-3.5pt) -- ++(0 pt,7.0pt);
\draw[color=blue] (-2.46,2.89) node {\large{$G$}};
\draw [color=blue] (0.18,1.62)-- ++(-3.5pt,0 pt) -- ++(7.0pt,0 pt) ++(-3.5pt,-3.5pt) -- ++(0 pt,7.0pt);
\draw[color=blue] (0.32,2.07) node {\large{$H$}};
\draw [color=blue] (-0.49985916889675985,4.437550622520741)-- ++(-3.5pt,0 pt) -- ++(7.0pt,0 pt) ++(-3.5pt,-3.5pt) -- ++(0 pt,7.0pt);
\draw[color=blue] (-0.36,4.89) node {\large{$I$}};
\draw [color=blue] (6.08,2.02)-- ++(-3.5pt,0 pt) -- ++(7.0pt,0 pt) ++(-3.5pt,-3.5pt) -- ++(0 pt,7.0pt);
\draw[color=blue] (6.02,2.55) node {\large{$J$}};
\draw [color=blue] (7.94,4.56)-- ++(-3.5pt,0 pt) -- ++(7.0pt,0 pt) ++(-3.5pt,-3.5pt) -- ++(0 pt,7.0pt);
\draw[color=blue] (8.08,5.01) node {\large{$L$}};
\draw [color=blue] (9.51142979696883,3.409267943951958)-- ++(-3.5pt,0 pt) -- ++(7.0pt,0 pt) ++(-3.5pt,-3.5pt) -- ++(0 pt,7.0pt);
\draw[color=blue] (9.66,3.85) node {\large{$K$}};
\end{scriptsize}
\end{tikzpicture}
\end{center}
\begin{enumerate}
\item Quelle est la nature du triangle $GHI$ ? Du triangle $DEF$ ? Du triangle $JKL$ ? Justifie tes réponses.

\medskip\noindent\grandscarreaux{\linewidth}{2.4cm}
\item Dans le triangle $DEF$, comment s'appelle le point $E$ ? Comment s'appelle le côté $[FD]$ ?

\medskip\noindent\grandscarreaux{\linewidth}{1.6cm}
\item Dans le triangle $JKL$, comment s'appelle le côté $[JK]$ ?

\medskip\noindent\grandscarreaux{\linewidth}{1.6cm}
\end{enumerate}

\medskip
\paragraph{Exercice 3}: \hfill \textbf{3 points}

\noindent Qui suis je ?

\begin{enumerate}
\item Je suis un quadrilatère ayant $4$ angles droits, mes côtés opposés sont de la même longueur. \medskip

\medskip\noindent\grandscarreaux{\linewidth}{0.8cm}
\item Je suis un quadrilatère, mes quatre côtés sont égaux cependant je n'ai pas d'angle droit. \medskip

\medskip\noindent\grandscarreaux{\linewidth}{0.8cm}
\item Je suis un quadrilatère, mes quatre côtés sont égaux et j'ai quatre angles droits.\medskip

\medskip\noindent\grandscarreaux{\linewidth}{0.8cm}
\end{enumerate}

\medskip
\paragraph{Exercice 4}: \hfill \textbf{4,5 points}

\medskip
\begin{enumerate}
\item Construire un triangle $ABC$ tel que $AB=4$cm, $AC=3$cm et $BC=3,5$cm.
\item Construire un triangle $DEF$ équilatéral tel que $DE=4$cm.
\item Construire un triangle $GHI$ rectangle en $H$ tel que $GH=4$cm et $GI=5$cm.
\end{enumerate}

\noindent \begin{minipage}{5.6cm}
\medskip\noindent\grandscarreaux{\linewidth}{5.6cm}
\end{minipage}
\hspace{2mm}
\begin{minipage}{5.6cm}
\medskip\noindent\grandscarreaux{\linewidth}{5.6cm}
\end{minipage}
\hspace{2mm}
\begin{minipage}{5.6cm}
\medskip\noindent\grandscarreaux{\linewidth}{5.6cm}
\end{minipage}


\medskip
\paragraph{Exercice 5}: \hfill \textbf{4,5 points}

\noindent Dans chaque cas, trace une figure à main levée puis réalise la figure en vraie grandeur.
\begin{enumerate}
\item Construire un rectangle $EFGH$ tel que $EF=4$cm et $FG=2$cm
\item Construire un losange $IJKL$ tel que $JL=3$cm et $IJ=4$cm
\item Construire un carré $MNOP$ tel que $MN=3$cm
\end{enumerate}

\noindent\begin{minipage}{5.6cm}
\medskip\noindent\grandscarreaux{\linewidth}{5.6cm}
\end{minipage}
\hspace{2mm}
\begin{minipage}{5.6cm}
\medskip\noindent\grandscarreaux{\linewidth}{5.6cm}
\end{minipage}
\hspace{2mm}
\begin{minipage}{5.6cm}
\medskip\noindent\grandscarreaux{\linewidth}{5.6cm}
\end{minipage}

\end{document}